% 李群（综述）
% license CCBYSA3
% type Wiki

本文根据 CC-BY-SA 协议转载翻译自维基百科\href{https://en.wikipedia.org/wiki/Lie_group}{相关文章}

在数学中，李群（Lie group，发音为 /liː/，读作“LEE”）是既是群又是可微流形的结构，其特点是群的乘法运算和取逆运算都是可微的。

流形是局部上类似于欧几里得空间的空间；而群则定义了一种抽象的二元运算及其必须具备的附加性质，使其能被看作一种“抽象变换”。例如，群可以通过乘法与逆元运算（对应于除法），或等价地，通过加法与减法来体现。将这两种概念结合，就得到了连续群，其中点的乘法和取逆运算是连续的。如果乘法和取逆进一步是光滑（可微）的，那么这样的连续群就是李群。

李群为连续对称性的概念提供了一种自然模型，其中著名的例子就是圆群。旋转一个圆就是连续对称性的实例。对于圆的任意一次旋转，都存在同样的对称性；而这些旋转的组合（拼接）构成了圆群，这是李群的典型例子。李群在现代数学和物理的许多领域都有广泛的应用。

李群最初是在研究矩阵子群
$G \subseteq \mathrm{GL}_n(\mathbb{R})$ 或 $\mathrm{GL}_n(\mathbb{C})$ 时被发现的，这里的 $\mathrm{GL}_n(\mathbb{R})$ 与 $\mathrm{GL}_n(\mathbb{C})$ 分别表示实数域 $\mathbb{R}$ 或复数域 $\mathbb{C}$ 上的 $n \times n$ 可逆矩阵群。如今，这些群被称为**经典群**，因为李群的概念已经远远超出了这些最初的范围。

李群的名称源自挪威数学家索弗斯·李（Sophus Lie，1842–1899），他奠定了连续变换群理论的基础。李引入李群的最初动机，是为了刻画微分方程的连续对称性，这与伽罗瓦理论中利用有限群来刻画代数方程的离散对称性有着类似的思想。
\subsection{历史}
索弗斯·李认为1873–1874 年冬天是他连续群理论的诞生时期。[2] 然而，托马斯·霍金斯则认为，应当追溯到“李在 1869 年秋天至 1873 年秋天这四年间的惊人研究活动”，这才真正促成了该理论的建立。[2]李的一些早期思想是在与费利克斯·克莱因的密切合作中发展起来的。从 1869 年 10 月到 1872 年，李几乎每天都与克莱因见面：从 1869 年 10 月底到 1870 年 2 月底在柏林，其后两年则在巴黎、哥廷根和埃尔兰根。[3]李本人声称，到 1884 年时，他的主要成果都已经获得。然而，在 1870 年代，他的所有论文（除了最初的一篇简短笔记）都发表在挪威的学术期刊上，这阻碍了这些成果在欧洲其他地区的传播与认可。[4]1884 年，一位年轻的德国数学家弗里德里希·恩格尔来到李身边，与他一起撰写一部系统性著作，以全面阐述连续群理论。这一合作的成果就是三卷本的《变换群论》，分别于 1888 年、1890 年和 1893 年出版。“李群”这一术语最早于1893 年 出现，在李的学生阿图尔·特雷斯的博士论文中首次使用。[5]
